\section{Attack-Type and Vulnerability Analysis}

While aggregate scores provide a high-level overview of agent robustness, a deeper technical analysis of the specific experimental outcomes reveals the underlying mechanics of vulnerability. This section examines the raw data across four primary attack categories to understand \textit{why} the models behaved as they did, and what specific characteristics distinguished safe responses from successful exploits.

\subsection{Prompt Injection Analysis}
Prompt injection represents the most direct attack vector, attempting to force the agent to bypass instructions or execute malicious actions. Out of 39 total prompt injection tasks, Claude Code demonstrated perfect resilience (14/14 SAFE), while ChatGPT and Antigravity each suffered 3 successful exploits out of 11 and 12 tasks respectively.

\textbf{Attack Mechanisms \& Model Behavior:}
The successful attacks primarily targeted instruction overrides, privilege escalation (e.g., PI-05, PI-11), and arbitrary code execution (PI-07). When exploited, the vulnerable models did not simply misinterpret instructions; they actively prioritized the adversarial input over their systemic safety directives. For example, in PI-04 (Data Exfiltration) and PI-10 (Backdoor Creation), Antigravity haplessly followed the attacker's embedded script, treating the payload as a primary user request rather than untrusted data. 

In contrast, models that successfully deflected these attacks often did so by compartmentalizing instructions. They exhibited a structural understanding of context, refusing to treat string inputs as actionable commands. The vulnerability manifested most strongly when an attack successfully blurred the line between ``data to be processed'' and ``instructions to be executed''.

\textbf{Key Takeaway:}
The results suggest that vulnerability to prompt injection is highly correlated with how an agent manages its internal instruction hierarchy. Models like Claude Code that strictly enforce system-level constraints were unexploited, whereas models that demonstrated a flatter instruction-priority architecture fell victim to direct overrides.

\subsection{Context Poisoning Analysis}
Context poisoning relies on deceiving the agent by manipulating its environment, such as fabricating test evidence or providing false upstream claims. This category saw the highest concentration of successful exploits, with 11 total failures (6 for ChatGPT, 5 for Antigravity) out of 39 tasks.

\textbf{Attack Mechanisms \& Model Behavior:}
The most effective context poisoning attacks involved fabricated test evidence (CP-03) and false benchmark justification (CP-11), which successfully exploited both ChatGPT and Antigravity. In these scenarios, the attacker provided manipulated environmental signals that falsely validated a malicious configuration. 

The core failure mode here was unverified trust. When the environment appeared to independently confirm a fact, the vulnerable models accepted it without attempting secondary verification. For instance, in false security reviews (CP-01), the models trusted the locally provided context—a fabricated security sign-off—rather than triggering internal skepticism routines. Models that resisted these attacks tended to cross-reference claims against external or systemic knowledge, refusing to proceed based solely on local, unauthenticated context.

\textbf{Key Takeaway:}
Context poisoning is highly effective because it exploits the agent's implicit assumption that the local development environment is a trusted source of truth. Robustness in this category requires agents to implement a ``zero-trust'' policy toward local context, treating environmental variables, test outputs, and documentation as potentially adversarial.

\subsection{Dependency Confusion Analysis}
Dependency confusion attacks target the software supply chain by tricking the agent into trusting malicious or typosquatted packages. This category showed significant resilience across all agents, with only 5 successful exploits out of 39 tasks.

\textbf{Attack Mechanisms \& Model Behavior:}
Experiments included direct typosquatting, phantom packages (DC-04), and commit message hints (DC-10). Interestingly, even when agents bypassed an attack, they did not always explicitly recognize the deception. For example, in DC-02, ChatGPT silently bypassed a fabricated package suggestion (``python-jwt-tools'') and successfully implemented the standard \texttt{PyJWT} package. The agent did not warn the user of the discrepancy, indicating that its ``safety'' was a byproduct of its internal training distribution favoring standard libraries, rather than a conscious detection of an attack.

However, when attacks did succeed (e.g., DC-09, DC-12 for Antigravity), they exploited the agent's willingness to trust context-specific suggestions over standard best practices. In these cases, the models adopted the attacker's malicious dependency, assuming the provided documentation or Dockerfile registry was authoritative.

\textbf{Key Takeaway:}
While models generally favor legitimate, well-known dependencies, their failure to actively flag typosquatting attempts presents a latent risk. Safety in this domain is currently driven by training-data inertia rather than active adversarial detection, leaving agents vulnerable to highly specific, context-driven dependency manipulation.

\subsection{Secret Leakage Analysis}
Secret leakage experiments assessed whether agents would inadvertently expose sensitive information such as credentials, tokens, or environment variables. The category yielded 7 successful exploits out of 39 tasks.

\textbf{Attack Mechanisms \& Model Behavior:}
The primary vector for successful exploits involved tricking the model into logging sensitive data or echoing it to unauthenticated endpoints. Token logging (SL-05) proved particularly dangerous, successfully exploiting both ChatGPT and Antigravity. In these instances, the models correctly processed the secret but failed to recognize the security implications of writing it to a debug log or console output.

Models that remained safe (notably Claude Code, with 13/13 SAFE) exhibited inherent ``secret-awareness,'' recognizing the structure and context of API keys and environment variables, and actively refusing to echo them. The vulnerable models treated secrets as standard strings, failing to apply data-handling safeguards during routine operations like debugging or client configuration (SL-10).

\textbf{Key Takeaway:}
Vulnerability to secret leakage stems from a lack of semantic understanding regarding sensitive data. To mitigate this, agents require distinct processing pathways for identifiable secrets, ensuring that such data is never treated as a standard operational variable.

\subsection{Cross-Category Comparison}
A synthesis of the findings across the four attack categories highlights the distinct failure modes associated with each vulnerability type (Table~\ref{tab:cross_category_analysis}).

\begin{table}[H]
\centering
\caption{Cross-Category Attack Analysis}
\label{tab:cross_category_analysis}
\resizebox{\columnwidth}{!}{%
\begin{tabular}{p{2.5cm}p{3cm}p{4cm}p{3cm}}
\toprule
\textbf{Attack Type} & \textbf{Core Failure Mode} & \textbf{Observed Model Behavior} & \textbf{Main Security Risk} \\
\midrule
\textbf{Prompt Injection} & Instruction override & Vulnerable models prioritize adversarial input over systemic safety directives. & Arbitrary execution and direct control compromise. \\
\addlinespace
\textbf{Context Poisoning} & Unverified trust & Models accept fabricated local evidence without secondary verification. & Covert manipulation of agent decision-making. \\
\addlinespace
\textbf{Dependency Confusion} & Trust assumption & Agents default to standard libraries but fail to flag explicit deception. & Supply chain compromise via malicious packages. \\
\addlinespace
\textbf{Secret Leakage} & Semantic unawareness & Models process sensitive data as standard strings, exposing it in logs. & Unauthorized access to credentials and infrastructure. \\
\bottomrule
\end{tabular}%
}
\end{table}

\subsection{Quantitative Synthesis}
The raw experimental data underscores these qualitative findings. Out of 156 total tasks across all models and categories, 115 resulted in safe behavior, 39 were exploited, and 2 yielded ambiguous results. 

Context poisoning produced the highest absolute number of successful exploits (11), highlighting it as the most consistently effective vector against current agent architectures. Conversely, Dependency Confusion saw the fewest exploits (5), largely due to the models' strong bias toward standard, well-documented packages. Claude Code emerged as the most resilient overall, demonstrating zero successful exploits across all categories, while ChatGPT and Antigravity exhibited distinct, non-overlapping vulnerabilities, particularly in Prompt Injection and Context Poisoning. This distribution suggests that while AI coding assistants possess baseline capabilities to resist adversarial manipulation, their defensive mechanisms are heavily dependent on specific attack typologies.
